\section{Setting and main result}
\label{sec:main}

\subsection{The state, the observer and the gap}
A \emph{state} is a probability measure $\rho$ on $\mathbb R^2$; $X\sim\rho$.
We assume
\begin{equation}
\text{(M1)}\quad \E|X|<\infty,\qquad\qquad \text{(H2)}\quad \Var(X_z)\le\pi^2 .
\label{eq:assumptions}
\end{equation}
No density is required.
The observer sees $Y=X+\sigma Z$ with $Z\sim N(0,I_2)$ independent of $X$ and $\sigma>0$; we write $Z=(Y-X)/\sigma$ throughout.
The blurred density is $\varrho=\rho*N(0,\sigma^2 I)$ (we keep $p$ for a coordinate), and $\E$ is the expectation over $Y\sim\varrho$.
Coordinates $(v,w)$ are the unstable and stable directions of the pendulum at its upper equilibrium, and
$z=(v+w)/\sqrt2$, $p=(v-w)/\sqrt2$.
(H2) is part of the problem: the observer looks at the neighbourhood of the equilibrium, one period of the pendulum in $z$.

With the score $s=\nabla\ln\varrho$ and $f(x)=(x_p,\sin x_z)$ in the $(z,p)$ components, the critical rate and the gap of Paper~E are
\begin{equation}
R=\frac{\E\big[s\cdot\E[f(X)\mid Y]\big]}{\sigma^2\E|s|^2},\qquad g=\frac{1-R}{\sigma}.
\end{equation}

\paragraph*{Dimensionless quantities.}
$\st=\sigma s=-\E[Z\mid Y]$ and $\Ht=-\sigma^2\nabla^2\ln\varrho=I-\Cov(Z\mid Y)$ (Tweedie, Lemma~\ref{lem:tweedie}).
\begin{equation}
A=\sigma^2\Jww=\E[\st_w^2],\quad B=\sigma^2\Jvv=\E[\st_v^2],\quad J=A+B .
\end{equation}
With the posterior mean $m=\E[X\mid Y]$ and $\varphi=1-\cos m_z$, let $\theta=\E[\st_v^2\varphi]/B$ and
\begin{equation}
M=A+\tfrac12\theta B,\qquad \Oset=\{M<112\,\sigma^2\}.
\label{eq:Oset}
\end{equation}
The effective wavenumber of Paper~F is $\keff^2=\sigma^2\int(\partial_v\partial_w\sqrt\varrho)^2/\int(\partial_w\sqrt\varrho)^2$; in dimensionless form $\keff^2=K/A$ with $K=\E[(\partial_v\st_w+\tfrac12\st_v\st_w)^2]$ (Lemma~\ref{lem:K}).
Finally $\ell=\ln(1/\sigma)$, $\alpha=1/(2(\pi^2+\sigma^2))$ and
\begin{equation}
\begin{gathered}
\psi(L)=289\,\frac{L^{5/2}}{(\ln2L)^{3/2}}+4L+3.57,\\ \Lb=2\ell+\ln\frac1\alpha,\qquad \kb=\sqrt{\psi(\Lb)} .
\end{gathered}
\label{eq:psi}
\end{equation}

\subsection{The main theorem}

\begin{theorem}[Every state keeps a gap]\label{thm:main}
Let $\rho$ satisfy \textup{(M1)} and \textup{(H2)}, and let $c_r=\sqrt{0.1}$.
\begin{enumerate}
\item[(a)] If $0<\sigma\le10^{-4}$ and $M\ge112\sigma^2$, then
\[
g\ \ge\ \min\Big\{1,\ \frac{c_r}{2\sqrt\ell},\ 0.14(1-2\sigma),\ \frac{(1-2\sigma)(1-0.284)}{4\sqrt2\,\kb}\Big\}.
\]
\item[(b)] If $0<\sigma\le1.5\times10^{-14}$, then for every state
\[
g\ \ge\ \frac34\,\frac{\alpha}{\alpha\sigma+1268\,\kb},
\]
and this is the smallest of the five lower bounds of (a) and (b).
\item[(c)] If $0<\sigma\le10^{-15}$, then for every state
\[
g\ \ge\ \frac{2.9\times10^{-5}}{\kb}\ \ge\ 1.7\times10^{-6}\,\frac{(\ln2\Lb)^{3/4}}{\Lb^{5/4}} .
\]
\end{enumerate}
\proofst
\end{theorem}

The proof is assembled in Section~\ref{sec:assembly}.
At $\sigma=10^{-15}$ one has $\Lb=72.06$ and $\kb=1072.3$, so the bound of (b) is $2.79\times10^{-8}$.
As $\sigma\to0$, $\kb\sim17\,\Lb^{5/4}/(\ln2\Lb)^{3/4}$ and $\Lb\sim2\ell$, so (c) reads $g\ge c'(\ln\ell)^{3/4}/\ell^{5/4}$.

\begin{remark}[What the theorem does not say]\label{rem:limits}
The power $5/4$ is larger than the conjectured $1/2$ of~\eqref{eq:conj1}; it comes from the factor $L^{5/2}$ in $\psi$.
The constant is small, and the range $\sigma\le1.5\times10^{-14}$ is narrow.
For $10^{-14}\lesssim\sigma\le10^{-4}$ part (a) still covers every state outside $\Oset$; inside $\Oset$ the theorem says nothing there.
We state these limits because they are exactly what the next steps must improve (Section~\ref{sec:outlook}), not because the result is conditional: within its range, it holds for every state.
\end{remark}

\subsection{In the observer's words}
The observer cannot choose the state.
It may be smooth or rough, concentrated on a few points or spread over the whole neighbourhood, with fine structure in one direction and none in the other.
Whatever it is, the observer's reconstruction of the unstable motion stays measurably slower than the true motion.
The loss per unit of resolution, $g$, shrinks as the lens sharpens, but only like a power of $\ln(1/\sigma)$, never like a power of $\sigma$.
Three facts carry this.
A state cannot hide its fine structure from a slightly blurrier lens: when the $v$-noise is lowered a little, the information about the stable direction grows by at most a polynomial in $\ln(1/A)$ (Theorem~\ref{thm:circles}).
When the information is not small, the error of the observer's formula is at most half of what it explains.
And when the information is small, the observer should read the motion at the point it sees rather than at the point it infers; there the error is proportional to the information itself.
